% 由 gewu 最小课题链路自动生成 —— 请勿手工修改，改模板后重跑 scripts/run_min_case.py
% 编译：xelatex main.tex（使用 ctexart 以支持中文）
\documentclass[11pt,a4paper,UTF8]{ctexart}
\usepackage{geometry}
\usepackage{booktabs}
\usepackage{amsmath}
\usepackage{hyperref}
\geometry{margin=2.5cm}
\title{多示例学习 + LIBS 热值预测\\[2pt]\large 合成数据上的端到端链路验证报告}
\author{格物 (Gewu) 最小课题自动链路}
\date{2026-11-09}
\begin{document}
\maketitle

\section*{摘要}
\noindent\textbf{本报告的全部数值来自合成光谱数据，用于验证「检索—复现—改进—成稿—投稿—沉淀」
链路可自主跑通，不代表真实煤样的精度结论。}
在 70 个 bag（共 611 条实例光谱）上，袋均谱 PLSR 基线 RMSE 为
$150.44$ kcal/kg；逐实例定标后池化（MIL）为 $122.79$ kcal/kg；进一步叠加强度上尾残差
（本方法）为 $73.03 \pm 10.17$ kcal/kg，相对基线降低 $51.46\%$。
逐折配对 $t$ 检验 $t(4)=-22.48$、$p<0.0001$，其中 5 折方向为改进。

\section{引言}
激光诱导击穿光谱（LIBS）用于煤质快速分析时，样品以「大颗粒堆积」形式被烧蚀：
一个煤样对应多条光谱（多示例），标签只在\textbf{样}级可得，单条光谱没有独立标签。
这正是多示例学习（MIL）的设定\cite{ilse2018mil}。
经典做法是把一个样品的多条光谱先平均成一条再回归\cite{li2014pls,yan2023coal}，
但平均会把颗粒尺度的差异抹平——少数粗颗粒主导信号，平均后这部分信息被稀释。

本文的贡献不在提出新网络，而是给出一条\textbf{可被第三方独立重跑}的最小链路：
数据可合成、模型只依赖 numpy/scikit-learn、每个数字都能回溯到 JSON Pointer。

\section{方法}
\subsection{数据生成机制}
合成数据把「成分」与「颗粒尺度」显式分离：
\begin{equation}
  \mathbf{x}_i = a_i \cdot \bigl[\, \mathbf{b} + \sum_k c_k \mathbf{p}_k \,\bigr] \cdot (1+\varepsilon_i) + \mathbf{n}_i,
  \qquad a_i = 1 + h z_i,
\end{equation}
其中 $c$ 为成分（逐实例可见），$h$ 为样品级颗粒异质性（\textbf{仅 bag 级可见}），
标签 $y = 3000 + 1200c_0 + 900c_1 + 600c_2 + 1400h + \mathcal{N}(0,20)$。
关键性质：$\overline{a_i} \to 1$ 使袋均谱几乎抹去 $h$，而上尾 $\max_i a_i$ 保留 $h$。

\subsection{三个对照臂}
\begin{itemize}
  \item \textbf{PLSR 基线}：袋均谱 $\to$ PLSR（$n_c=6$）；
  \item \textbf{MIL 直接}：逐实例谱形定标 $\mathbf{x}_i/\overline{\mathbf{x}_i}$ $\to$ 实例 PLS $\to$
        LSE 池化（$\beta$ 由内层 CV 选，$\beta=0$ 即均值）；
  \item \textbf{本方法}：MIL 直接 + 以强度池化统计量（max/p90/std/极差/实例数，共 7 维）回归残差的 Ridge。
\end{itemize}

\section{实验设置}
5 折交叉验证，\textbf{按 bag 切分}（同一 bag 的实例不会跨折，否则泄漏）。
$\beta$ 与 Ridge 的 $\alpha$ 只在内层 4 折的折外预测上选择；残差特征同样取折外预测，
避免自拟合乐观偏差。随机种子固定，重复运行结果逐位一致。

\section{结果}
\begin{table}[h]
\centering
\caption{5 折交叉验证（RMSE，kcal/kg）}
\begin{tabular}{lccc}
\toprule
臂 & 均值 & 总体标准差 & 逐折 \\
\midrule
PLSR 基线 & 150.44 & 13.15 & 169.63 / 134.72 / 158.11 / 136.84 / 152.88 \\
MIL 直接 & 122.79 & 12.24 & 141.31 / 109.24 / 118.55 / 112.30 / 132.54 \\
本方法 & 73.03 & 10.17 & 91.87 / 69.17 / 72.40 / 61.18 / 70.52 \\
\bottomrule
\end{tabular}
\end{table}

本方法 5 折平均 $R^2=0.911$。逐折配对 $t$ 检验（本方法 $-$ 基线）：
$t(4)=-22.48$，$p<0.0001$，95\% 置信区间 $[-86.97, -67.85]$；
Wilcoxon 检验 $p=0.0625$（$n=5$ 时双侧最小可达 $0.0625$，本身不显著，两种检验结论一并报告）。

\section{消融与讨论}
在三臂对比中，去掉残差（MIL 直接）比基线低 $18.38\%$，再加残差再低 $33.08\%$。
这说明本数据上的增益主要来自两点：\textbf{逐实例定标}把尺度与成分解耦；
\textbf{强度上尾统计量}回收了被袋均抹掉的 $h$；残差模型只吃基线看不到的量，刻意不放入与 MIL 输出共线的预测统计量（实测会过拟合）。
内层 CV 选出的 $\beta$ 为 $0, 0.25, 0.25, 0.25, 0.25$——若 $\beta=0$ 占多数，应如实解读为
「该数据上均值池化已足够，增益不来自高阶池化」，而不是宣称 LSE 有效。

\paragraph{适用边界。}合成数据不含基体效应、自吸收与仪器漂移，绝对误差量级不可外推到真实煤样；
本报告证明的是\textbf{机制与链路}，真实精度以存档课题的实测指标为准。

\section{结论}
在零外部数据依赖的前提下，七阶段链路可从零跑出可复现的数值、可回溯的断言与可编译的成稿，
为「成果可复现性」提供了第三方可独立执行的证据。

\begin{thebibliography}{9}
\bibitem{li2014pls} Y. Li, et al. A model combining spectrum standardization and dominant factor based partial least square method for carbon analysis in coal by LIBS. \textit{Spectrochimica Acta Part B}, 2014. DOI: 10.1016/j.sab.2014.06.017
\bibitem{ilse2018mil} M. Ilse, J. Tomczak, M. Welling. Attention-based Deep Multiple Instance Learning. \textit{ICML}, 2018. arXiv:1802.04712
\bibitem{yan2023coal} Yan, et al. Rapid Classification and Quantification of Coal by Using Laser-Induced Breakdown Spectroscopy and Machine Learning. \textit{Applied Sciences}, 2023. DOI: 10.3390/app13148158
\end{thebibliography}
\end{document}
